\documentclass{subfiles}
\begin{document}
    We now discuss an algorithm that can be used to solve, when it works, the DLP.
        Hence, this algorithm is a good tool for cracking ElGamal-based cryptosystems.

    The idea is the following: let \(S\) be the set of elements we are working with,
        and consider a function \(f : S \to S\) that randomly distributes the elements in \(S\).
        Since for our purpouse \(S\) is finite, by applying \(f\) repetedly 
        it will happen that 
        \[
            f(x_{i}) = f(x_{j}), i \ne j.
        \]
    The only problem is to determine such a function.

    In our context of the ElGamal algorithm, let \(\gf{p}\) be the choosen field 
        and let \(g\) be a primitive root. Then, as prosed in (\cite{pollard1996}),
        we can use the function
        \[
            f(x) = \begin{cases}
                gx, \text{if } 0 \le x < \sfrac{p}{3}; \\ 
                x^{2}, \text{if } \sfrac{p}{3} \le x < \sfrac{2p}{3}; \\ 
                ax, \text{if } \sfrac{2p}{3} \le x < p;
            \end{cases}
        \]
        where \(x\) is reduced modulo \(p\). We can assume that \(x_{0} = 1\) and
        by applying \(f\) we get simply 
        \[
            g^{\alpha_{i}}a^{\beta_{i}}
        \]
        where the values of \(\alpha_{i}, \beta_{i}\) are given by 
        \[\alpha_{i} = \begin{cases}
                \alpha_{i - 1} + 1, \text{if } 0 \le x < \sfrac{p}{3}; \\ 
                2\alpha_{i - 1}, \text{if } \sfrac{p}{3} \le x < \sfrac{2p}{3}; \\ 
                \alpha_{i - 1}, \text{if } \sfrac{2p}{3} \le x < p;
            \end{cases}\]


        \[\beta_{i} = \begin{cases}
            \beta_{i - 1}, \text{if } 0 \le x < \sfrac{p}{3}; \\ 
            2\beta_{i - 1}, \text{if } \sfrac{p}{3} \le x < \sfrac{2p}{3}; \\ 
            \beta_{i - 1} + 1, \text{if } \sfrac{2p}{3} \le x < p.
        \end{cases}.\]

        Then, if we compute \(y_{i} = f(f(y_{i - 1}))\), with \(y_{0} = x_{0}\)
            and check when \(y_{i} = x_{i}\) we get something of the form 
            \[
                g^{\alpha_{i}}a^{\beta_{i}} = g^{\gamma_{i}}a^{\delta_{i}}
            \]
            where \(\gamma_{i}, \delta_{i}\) are defined by the same relation of 
            \(\alpha \text{ and } \beta\).
            Hence, we have that 
            \[
                g^{\alpha_{i} - \gamma_{i}} = a^{\beta_{i} - \delta_{i}}.
            \]
            Putting \(u = \alpha_{i} - \gamma_{i}\) and \(v = \beta_{i} - \gamma_{i}\)
            and considering the logarithm, we get 
            \[
                v \log_{g} a \equiv u \mod{(p - 1)}.
            \]
            Recall that \(p\) is known, therefore, if \(\gcd{v, p - 1} = 1\)
            we can multiply by \(v^{-1}\) and easely compute the logarithm.
            Viceversa, if \(\gcd{v, p - 1} \ge 2\), we can apply the Eucledian algorithm
            to find \(s\) such that \(su \equiv d \mod{(p - 1)}\). Multiplying 
            both sides by \(s\), we get 
            \[
                d \log_{g}{a} \equiv su \mod{(p - 1)},
            \]
            whose solutions is the set 
            \[
                \set*{\frac{w}{d} + k \frac{p - 1}{d}}[0 \le k < d].
            \]
\end{document}
